\documentclass{ieeeaccess}
\usepackage{cite}
\usepackage{amsmath,amssymb,amsfonts}
\usepackage{algorithm}
\usepackage{algpseudocode}
\usepackage{graphicx}
\usepackage{textcomp}
\usepackage{xcolor}
\usepackage{hyperref}
\usepackage{url}
\usepackage{booktabs}
\usepackage{multirow}
\usepackage{float}
\usepackage{lipsum} % For generating filler text if needed to expand sections, though we will provide real text.

% Compatibility guard for IEEE Access class versions whose caption code
% references this dimension register without declaring it.
\makeatletter
\@ifundefined{xfigwd}{\newlength{\xfigwd}}{}
\makeatother

\def\BibTeX{{\rm B\kern-.05em{\sc i\kern-.025em b}\kern-.08em
    T\kern-.1667em\lower.7ex\hbox{E}\kern-.125emX}}
\begin{document}
\history{Date of publication xxxx 00, 0000, date of current version xxxx 00, 0000.}
\doi{10.1109/ACCESS.2017.DOI}

\title{A Learned Adaptive Inference System for Real-Time YOLOv8 Object Detection on Edge Devices with Pothole Detection as an Application}

\author{
\uppercase{Vishal K. Bhat}\authorrefmark{1},
\uppercase{Vaibhav P. Rodappanavar}\authorrefmark{1},
\uppercase{Naren Sridharan}\authorrefmark{2},
\uppercase{Hardik Paneru}\authorrefmark{2}, and
\uppercase{Rajeswara Rao K V S}\authorrefmark{3}}

\address[1]{Department of Computer Science, RV College of Engineering, Bengaluru, Karnataka 560059, India (e-mail: vishalkbhat.cs23@rvce.edu.in; vaibhavp.cs23@rvce.edu.in)}
\address[2]{Department of Electronics and Communication Engineering, RV College of Engineering, Bengaluru, Karnataka 560059, India (e-mail: narensridharan.ec23@rvce.edu.in; hardikpaneru.ec23@rvce.edu.in)}
\address[3]{Department of Industrial Engineering \& Management, RV College of Engineering, Bengaluru, Karnataka 560059, India (e-mail: rajeswararaokvs@rvce.edu.in)}

\tfootnote{This work was supported by the Department of Computer Science at RV College of Engineering.}

\markboth
{Bhat \headeretal: A Learned Adaptive Inference System for Real-Time YOLOv8 Pothole Detection}
{Bhat \headeretal: A Learned Adaptive Inference System for Real-Time YOLOv8 Pothole Detection}

\corresp{Corresponding author: Rajeswara Rao K V S (e-mail: rajeswararaokvs@rvce.edu.in).}

\begin{abstract}
Deploying real-time object detection on resource-constrained edge hardware such as the Raspberry Pi presents a fundamental challenge: sustained inference at a fixed configuration inevitably drives the processor into thermal throttling, causing frame rates to collapse and detections to be missed at safety-critical moments. Static pipelines offer no mechanism to reconcile the gap between peak computational demand and the hardware's available thermal headroom. This paper presents a Learned Adaptive Inference System that treats edge deployment as a continuous hardware-aware control problem. The system monitors five live telemetry signals—CPU utilization, RAM consumption, core temperature, instantaneous frame rate, and per-frame latency—and uses these to drive a lightweight neural controller that continuously outputs the optimal input resolution and frame-skip interval for a YOLOv8n pothole detection pipeline.

The system is evaluated on a real-world pothole detection task using a YOLOv8n model deployed on a Raspberry Pi 5, where the thermal constraints of the hardware make adaptive inference particularly consequential.

The learned controller reduces average input resolution from 640\,px to 400\,px under load, yielding a 61.0\% reduction in pixel-level compute with only a 1.8\% drop in F1-score. Counter-intuitively, operating at reduced resolution increases detection precision by 1.2\% by suppressing false positives caused by road texture and shadow artifacts. These results demonstrate that offline reinforcement learning provides a principled, low-overhead path to thermally stable, continuous pothole detection on edge devices without any hardware modification or model retraining.
\end{abstract}

\begin{keywords}
Adaptive Inference, YOLOv8, Edge Computing, Offline Reinforcement Learning, Thermal Management, Pothole Detection, Raspberry Pi, Dynamic Resolution Scaling, Frame Skipping, Telemetry-Driven Control, Resource-Constrained Inference, IoT Edge Deployment.
\end{keywords}

\titlepgskip=-15pt

\maketitle

\section{Introduction}
\label{sec:introduction}
\PARstart{T}{he} deployment of Internet of Things (IoT) devices and autonomous systems has grown significantly, creating a critical need for performing complex computer vision tasks directly on resource-constrained edge devices. Among these tasks, real-time object detection plays an important role in applications requiring timely environmental perception and decision-making. The YOLO family of architectures, particularly YOLOv8, has gained significant attention for its favorable balance between detection accuracy and inference speed, making it well suited for real-time object detection on edge platforms. However, deploying deep learning-based object detection models on edge devices presents significant challenges due to inherent hardware limitations, including constrained computational resources, memory capacity, power availability, and thermal constraints. These challenges become particularly relevant in continuous real-time applications such as pothole detection, where the detection model must process a persistent stream of visual data while operating within the limited resources of the edge device.

Edge devices such as the Raspberry Pi 5 operate under substantially tighter computational, memory, power, and thermal constraints than conventional desktop or server-class systems. These limitations become particularly significant when deep learning-based object detection is performed continuously on live video streams. A YOLOv8 model executed at a fixed input resolution and frame-processing rate imposes a relatively constant computational workload, which can lead to sustained processor utilization and increasing thermal load during prolonged operation. As the device approaches its thermal limits, frequency throttling can reduce available computational performance, resulting in increased per-frame latency, reduced frame rates, and missed detections. Consequently, a static inference configuration cannot reliably maintain real-time performance under varying hardware conditions, motivating the need for an inference mechanism that dynamically adapts its computational workload to the available resources and thermal headroom.

In practical road-monitoring applications, pothole detection systems may need to operate continuously on edge devices deployed on vehicles for extended periods and under varying environmental and operating conditions. A YOLO-based detection model running continuously on such hardware cannot assume that the available computational and thermal resources will remain at their maximum capacity throughout the entire deployment. Sustained inference at a fixed configuration can increase processor utilization and thermal load, potentially leading to performance degradation through increased latency, reduced frame rates, and thermal throttling. Consequently, maintaining reliable real-time pothole detection requires the inference workload to adapt to the current state of the edge device rather than relying on a fixed configuration. Dynamically adjusting the computational workload according to live hardware and inference conditions can therefore help sustain real-time operation and extend the effective operating period of the detection system under resource-constrained conditions.

\subsection{The Limitation of Static Pipelines}
Current deployment strategies generally treat inference pipelines as static configurations. An engineer selects a fixed input resolution, such as $640 \times 640$, and frames from the camera sensor are then processed sequentially. While this approach theoretically yields consistent inference behavior under ideal conditions, it fails to account for the dynamic physical state of the hardware.

The static paradigm assumes that the edge device can consistently dissipate heat and sustain peak processing speeds. However, when thermal throttling occurs, the software continues demanding a fixed computational workload while the hardware's available computational capacity diminishes. This mismatch can result in reduced frame rates, increased latency, missed detections, and, in the worst case, system failure. A robust edge AI system must monitor its operational state and dynamically adjust its computational demands to preserve overall stability and responsiveness.

\subsection{Proposed Solution: Learned Adaptive Inference}
To address these challenges, a Learned Adaptive Inference System is introduced. This system transforms the standard YOLOv8 inference pipeline into an adaptive process that continuously responds to the live hardware state. The pipeline exposes two tunable control parameters at runtime: the input resolution $imgsz \in [256, 640]$ (constrained to multiples of 32) and the frame-skip probability $p_{skip} \in [0, 1]$, which determines the proportion of incoming frames that bypass the forward pass entirely. Together, these parameters form a continuous action space that allows fine-grained, real-time trade-offs between detection fidelity and computational demand.

A dedicated monitoring module samples five hardware telemetry signals asynchronously at 2\,Hz—CPU utilization, RAM utilization, core temperature, smoothed inference frame rate, and per-frame latency—maintaining a thread-safe state vector $s_t$ accessible to the inference loop in $\mathcal{O}(1)$ time. A lightweight neural Controller reads $s_t$ at each decision step and emits a continuous action pair $(imgsz_t,\; p_{skip,t})$, dynamically balancing detection quality against the available thermal and compute headroom of the edge device.

The resource management problem is formulated as a Markov Decision Process (MDP), enabling the system to be developed and evaluated through a rigorous offline reinforcement learning workflow:
\begin{enumerate}
    \item \textbf{Baseline Profiling}: Characterizing the thermal failure modes of a fixed $640 \times 640$ static pipeline under sustained load, establishing the performance upper bound and the conditions under which degradation occurs.
    \item \textbf{Continuous Action Space Exploration}: An exploration controller randomly samples continuous $(imgsz, p_{skip})$ pairs across the full parameter range, executing inference and logging state-action-reward-next-state transition tuples to build the offline RL dataset.
    \item \textbf{Heuristic Control (Baseline)}: A rule-based controller that selects from three discrete configurations—$M_0\!=\!(640,0)$, $M_1\!=\!(480,1)$, $M_2\!=\!(320,2)$—using CPU load, RAM utilization, and temperature thresholds. This serves as a strong interpretable baseline against which the learned policy is evaluated.
    \item \textbf{Offline RL via Filtered Behavioral Cloning}: The collected transitions are filtered to retain the top 30\% by reward. A Continuous Actor network is then trained offline via supervised regression on these high-quality transitions, learning to map $s_t$ directly to continuous $(imgsz, p_{skip})$ actions. This offline reinforcement learning approach produces a policy that has internalized the hardware's thermal dynamics without requiring any online interaction during training.
\end{enumerate}

\subsection{Structure of the Paper}
The remainder of this paper is organized as follows: Section \ref{sec:related_work} reviews existing research on edge object detection, thermal management, and dynamic inference. Section \ref{sec:architecture} describes the proposed system architecture, including hardware monitoring and inference execution layers. The step-by-step methodology, progressing from baseline profiling to reinforcement learning, is detailed in Section \ref{sec:methodology}. Experimental results are presented in Section \ref{sec:experiments}, followed by a discussion of the findings and limitations in Section \ref{sec:discussion}. Finally, Section \ref{sec:conclusion} concludes the paper and suggests directions for future research.

\section{Related Work}
\label{sec:related_work}
The fields of edge computing, deep learning, and adaptive resource management have evolved rapidly, with recent research focusing on improving detection accuracy while reducing computational, energy, and thermal costs. The present work builds upon these research directions while addressing the specific challenge of maintaining stable real-time object detection performance under prolonged thermal stress on resource-constrained edge devices.

\subsection{Pothole Detection on Edge Devices}

Deep learning-based object detection has become an effective approach for automated road-damage and pothole detection. Recent datasets and studies have established the use of deep learning for road-surface assessment, while improved YOLO-based architectures have demonstrated promising accuracy for pothole detection in diverse road environments \cite{b1,b2,b5,b6}. In particular, recent YOLOv8-based approaches have investigated improved feature extraction, segmentation, and detection performance for complex road conditions \cite{b3,b11,b12}. Lightweight road-distress detection models have also been explored for deployment on mobile and resource-constrained platforms \cite{b6}.

The suitability of these models for edge deployment has been further demonstrated through real-time object detection studies on embedded hardware \cite{b13,b14}. However, most existing approaches primarily optimize the model architecture or inference implementation for computational efficiency. Such static optimizations do not directly address the performance degradation that can occur when an edge device performs continuous inference and accumulates heat over extended periods. The proposed system instead introduces a runtime pipeline-level adaptation mechanism that dynamically adjusts input resolution and processing frequency according to the device's thermal state, allowing the same YOLOv8n model to operate under changing thermal conditions without retraining or replacing the model.

\subsection{Dynamic Neural Networks and Adaptive Inference}

Dynamic Neural Networks provide a prominent approach for reducing inference computation by adapting the amount of processing performed for each input. BranchyNet introduced early-exit branches that allow samples with sufficiently confident predictions to terminate inference before reaching the full network depth, thereby reducing computational cost \cite{b23}. Similarly, dynamic channel pruning selectively suppresses computational channels to reduce the workload associated with neural network inference \cite{b21}.

More recent work has extended adaptive inference to object detection. AdaDet introduces an early-exit mechanism specifically for object detection and uses detection uncertainty to determine whether an input can terminate computation early \cite{b17}. Other approaches have investigated adaptive inference directly on edge devices, dynamically adjusting inference behavior according to available computational resources \cite{b18}. Energy-aware approaches have further combined neural-network optimization with hardware-level frequency scaling to reduce energy consumption while satisfying latency constraints \cite{b8}.

Although these techniques can provide substantial computational savings, many require architectural modifications, additional exit branches, specialized policies, or retraining. Their computational savings may also depend on the complexity of individual inputs. In contrast, the proposed approach operates at the inference-pipeline level: input resolution and frame-processing frequency are dynamically adjusted while retaining a single, unmodified YOLOv8n model. This allows the adaptation mechanism to be applied without retraining and provides direct control over the workload imposed on the edge device.

\subsection{Thermal Management in Edge Computing}

Thermal management is an important consideration for embedded platforms executing computationally intensive workloads. Conventional operating-system-level thermal management can rely on frequency scaling when device temperatures approach critical limits, which can lead to reduced computational performance and violations of real-time requirements. Maity \emph{et al.} investigated thermal management for heterogeneous CPU-GPU embedded platforms and demonstrated a predictive control strategy that considers workload scheduling and frequency tuning to reduce peak temperature while maintaining real-time requirements \cite{b19}.

Recent energy-efficient inference research has also explored the interaction between neural-network computation and hardware power management. Dai \emph{et al.} proposed a software--hardware co-design approach that combines neural-network optimization, early exits, and frequency scaling to reduce energy consumption while satisfying latency constraints \cite{b8}. These studies demonstrate the importance of considering both computational workload and hardware operating conditions when deploying deep-learning workloads on edge devices.

However, the aforementioned approaches primarily focus on hardware scheduling, frequency management, or energy optimization. The proposed system instead establishes a direct feedback loop between hardware temperature telemetry and the computer-vision inference pipeline. When the measured temperature increases, the controller reduces the computational workload by modifying inference resolution and frame-processing frequency before severe thermal throttling occurs. This application-level strategy is designed specifically for continuous video inference and aims to maintain stable detection performance while keeping the device within a safe thermal operating range.

\subsection{Reinforcement Learning and Edge-Cloud Resource Optimization}

Reinforcement Learning (RL) has increasingly been investigated for resource
management in edge-computing environments. Recent surveys have examined
RL-based approaches for mobile edge computing, including resource allocation,
task scheduling, and computation offloading \cite{b7}. For example,
dependency-aware reinforcement learning has been used to formulate dynamic
task-offloading decisions as an optimization problem involving changing
computational and communication resources \cite{b9}.

Deterministic task-offloading approaches have also been investigated for the
IoT-edge-cloud continuum. Aghababaiyan \emph{et al.} proposed a deterministic
task-offloading and resource-allocation scheme that jointly considers
communication and computing resources across IoT devices, edge servers, and
cloud infrastructure \cite{b25}. Their approach prioritizes completing tasks
within their deadlines rather than minimizing the execution latency of each
individual task, thereby providing deterministic service levels under
resource-constrained conditions. Such approaches demonstrate that
computation can be dynamically distributed across the IoT-edge-cloud
continuum when communication and remote computational resources are available.

Adaptive inference has also been explored as a means of dynamically modifying
the computational resources allocated to deep neural network workloads.
Edge-device adaptive inference approaches demonstrate that inference behavior
can be modified according to device conditions and workload requirements
\cite{b18}. These studies motivate the use of feedback-driven decision-making
for managing computational workloads on resource-constrained edge devices.

The present work differs from task-offloading approaches in its control
objective and deployment architecture. Rather than distributing inference
tasks across remote computing resources or selecting among multiple neural
network architectures, the proposed system adapts the execution parameters of
a single YOLOv8n detector locally. Specifically, the learned controller
adjusts the input resolution and frame-skipping probability according to
real-time hardware and inference telemetry. This provides workload adaptation
without requiring network connectivity or transferring inference workloads to
external computing infrastructure. The heuristic controller is retained as an
interpretable baseline, while the final learned controller provides
continuous adaptation of the inference parameters.

\section{System Architecture}
\label{sec:architecture}
The system architecture is designed to decouple hardware monitoring, the decision-making controller, and the computationally intensive inference execution. This separation ensures that the overhead of adaptivity does not impede the detection pipeline itself.

\begin{figure*}[t!]
    \centering
    \includegraphics[width=0.85\textwidth]{placeholder_architecture.png}
    \caption{Comprehensive System Architecture showing the asynchronous interaction between the Hardware Monitoring Thread, the State Space Aggregator, the Adaptive Controller (Heuristic/Learned), and the YOLOv8 Inference Pipeline.}
    \label{fig:system_architecture_detailed}
\end{figure*}

\subsection{Hardware Abstraction and Telemetry Layer}

At the foundation of the system lies a Hardware Abstraction Layer (HAL) that
continuously monitors the operational state of the target edge device. The
HAL runs in a dedicated background thread, allowing telemetry to be collected
without blocking the real-time inference loop. In this work, the proposed
system is implemented and evaluated on a Raspberry Pi 5 as the target
resource-constrained edge platform.

The monitored hardware metrics comprise CPU utilization ($U_{cpu}$), RAM
utilization ($U_{ram}$), and core temperature ($T_{core}$). These are combined
with two inference-derived metrics: the smoothed frame rate ($FPS_{avg}$) and
per-frame latency ($L_t$). Together, these five signals form the complete
state vector $s_t$ supplied to the adaptive controller:

\begin{equation}
s_t =
[U_{cpu},\, U_{ram},\, T_{core},\, FPS_{avg},\, L_t].
\end{equation}

The HAL samples these signals asynchronously at a fixed interval of 500~ms
(2~Hz) and stores the latest telemetry snapshot in a thread-safe in-memory
structure. The inference pipeline can therefore access the current state
without waiting for a telemetry update, introducing negligible monitoring
overhead to the detection process.

\subsection{The Adaptive Controller}

The adaptive controller serves as the decision-making component of the
system, receiving the latest telemetry state from the Hardware Abstraction
Layer. At each control step, the controller processes the five-dimensional
state vector
$s_t = [U_{cpu}, U_{ram}, T_{core}, FPS_{avg}, L_t]$
and produces a continuous action consisting of the raw image-resolution and
frame-skipping parameters.

The learned controller is implemented as a lightweight Continuous Actor
network with a $5 \rightarrow 128 \rightarrow 128 \rightarrow 2$ multilayer
perceptron (MLP). ReLU activations are used in the hidden layers, while a
Tanh activation at the output constrains the two raw action values to the
range $[-1,1]$. These outputs are subsequently mapped to the operational
inference parameters: the image resolution $imgsz_t$ is mapped to the range
$[256,640]$ and constrained to multiples of 32, while the frame-skipping
probability $p_{skip,t}$ is mapped to $[0,1]$.

The controller is trained offline using reward-filtered behavioral cloning.
During the data-collection phase, continuous actions are explored on the
target edge platform and the resulting state-action-reward transitions are
recorded. The highest-reward transitions are then used as demonstrations for
training the Continuous Actor to learn a direct mapping from the observed
hardware and inference state to an appropriate inference configuration.
Consequently, no online learning or parameter updates are performed during
deployment.

At runtime, the resulting action pair $(imgsz_t,p_{skip,t})$ is passed to the
inference pipeline. This allows the computational workload to be continuously
adjusted according to the current operating state of the edge device rather
than relying on a fixed inference configuration.

\subsection{The YOLOv8n Inference Pipeline}

The visual processing stage is implemented using the Ultralytics YOLOv8n
detector. The inference pipeline receives video frames together with the
current control parameters $(imgsz_t, p_{skip,t})$ generated by the adaptive
controller. For frames selected for processing, the input image is resized to
the controller-selected resolution $imgsz_t$ and passed through the YOLOv8n
forward pass to produce object detections. The resulting predictions are
subsequently post-processed to obtain the final bounding-box detections.

The frame-skipping mechanism is probabilistic, with the controller output
$p_{skip,t}$ defining the probability that an incoming frame is bypassed.
Thus,
\begin{equation}
P(\mathrm{skip}_t) = p_{skip,t},
\qquad
P(\mathrm{process}_t) = 1-p_{skip,t}.
\end{equation}
Frames selected for skipping bypass the inference pipeline, avoiding the
computational cost of preprocessing, forward inference, and post-processing.
Frames selected for processing are resized according to the current
$imgsz_t$ configuration and passed to the YOLOv8n detector.

The inference pipeline also provides feedback to the adaptive controller in
the form of the smoothed frame rate $FPS_{avg}$ and per-frame latency $L_t$.
These inference-derived measurements are incorporated into the telemetry
state used at the subsequent control step. The resulting closed-loop
architecture allows the computational workload of the YOLOv8n pipeline to be
continuously adjusted according to the current hardware and inference
conditions rather than relying on a fixed inference configuration.

\section{Methodology}

\label{sec:methodology}

The adaptive system was developed through a rigorous, multi-phase sequence.
This systematic progression ensures that the final learned policy is
evaluated against well-understood baselines rather than arbitrary
benchmarks.

\subsection{Phase 0: Baseline Static Inference}

The objective of Phase 0 is to quantify the performance limitations of a
fully static inference pipeline on the target edge hardware, providing an
empirical baseline that motivates the need for adaptive control.

In this phase, the system operates as a standard static pipeline. YOLOv8n is
configured with a fixed input resolution of $640 \times 640$ pixels and FP32
precision, with no frame skipping. Every incoming frame is processed
unconditionally, resulting in a consistently high computational workload.

During execution, the system records the timestamp, frame rate, per-frame
latency, CPU utilization, RAM utilization, core temperature, number of
detections, average detection confidence, and cumulative frame-drop rate for
each processed frame.

On the Raspberry Pi~5 target platform, the static pipeline demonstrates
limited real-time performance. Across the evaluation run, the achieved frame
rate ranges from $1.8$ to $4.01$~FPS, while per-frame latency ranges from
approximately $245$ to $665$~ms. The cumulative frame-drop rate remains high,
ranging from $85.73\%$ to $96.88\%$ and settling at approximately $87.04\%$
toward the end of the run. Core temperature increases from approximately
$50^\circ$C to a maximum of $60.6^\circ$C, while RAM utilization remains
relatively stable around $30\%$. The observed temperature range provides no
evidence of thermal throttling during this particular run, indicating that
the primary limitation of the static configuration is the computational
throughput required for full-resolution YOLOv8n inference on the Raspberry
Pi~5.

The resulting behavior is illustrated in Fig.~\ref{fig:phase0_graphs}. The
figure shows the low sustained frame rate and high inference latency, together
with substantial frame loss throughout the run. These results demonstrate
that a fixed $640 \times 640$ inference configuration is unsuitable for
sustained real-time operation under the tested edge-device conditions.
Reducing the computational workload through lower input resolution, frame
skipping, or a combination of both therefore provides a basis for the
adaptive inference strategy developed in the subsequent phases.

\begin{figure*}[t!]
    \centering
    \includegraphics[width=0.85\textwidth]{logs_rpi5_overview.png}
    \caption{Phase~0 static baseline on the Raspberry Pi~5 using $640 \times
    640$ input resolution, FP32 precision, and no frame skipping. \textit{Top}:
    FPS and per-frame latency over the evaluation run. \textit{Middle}: CPU
    utilization, RAM utilization, and core temperature. \textit{Bottom}:
    detection count, average detection confidence, and cumulative frame-drop
    rate. The static configuration exhibits low sustained frame rates,
    substantial inference latency, and a high frame-drop rate, demonstrating
    the limitations of fixed full-resolution inference under the resource
    constraints of the target edge platform.}
    \label{fig:phase0_graphs}
\end{figure*}


\subsection{Phase 1: Action Space Definition and Controlled Adaptation}

Phase~0 establishes that a static $640 \times 640$ pipeline imposes a
substantial computational workload on the Raspberry Pi~5. Phase~1 therefore
defines the controllable parameters of the inference pipeline and
experimentally characterizes their effect on computational load and inference
performance.

The action space is based on two independent workload-control mechanisms:
input resolution and inference frequency. The input resolution $imgsz$
determines the spatial dimensions of the tensor presented to YOLOv8n, while
the frame-skipping parameter controls how frequently incoming frames are
processed.

\subsubsection*{Rationale for Resolution Scaling}

The computational cost of convolutional processing depends strongly on the
spatial dimensions of the input tensor. Reducing the input resolution
therefore provides a direct mechanism for reducing the computational workload
without modifying the YOLOv8n model architecture or its learned parameters.
In particular, reducing the input from $640 \times 640$ to
$320 \times 320$ reduces the input spatial area by $75\%$, providing a
substantial reduction in the amount of spatial computation performed by the
network.

\subsubsection*{Rationale for Frame Skipping}

Resolution scaling reduces the computational cost of each processed frame,
whereas frame skipping reduces the number of frames on which inference is
performed. These mechanisms therefore provide complementary forms of
workload control. For the heuristic baseline, three discrete configurations
are defined:

\begin{itemize}
    \item $M_0=(640,0)$: full-resolution inference with no frame skipping.
    
    \item $M_1=(480,1)$: reduced-resolution inference with every alternate
    frame bypassed.
    
    \item $M_2=(320,2)$: reduced-resolution inference with two out of every
    three frames bypassed.
\end{itemize}

These configurations span progressively lower computational workloads and are
used to characterize the resolution--throughput trade-off on the target
hardware.

During Phase~1, the feedback controller is disabled and each configuration is
applied for a fixed duration. The resulting frame rate, latency, CPU
utilization, and temperature are recorded to quantify the computational and
thermal characteristics of each operating point. Detection count and average
detection confidence are also recorded to assess the effect of workload
reduction on detection behavior. These measurements provide the empirical
basis for the heuristic thresholds introduced in the subsequent phase.

\subsubsection*{Phase~1 Observations and Trade-off Analysis}

The experimental results demonstrate a clear trade-off between computational
workload and detection performance, as shown in Fig.~\ref{fig:phase1_mode_summary}.
As the operating configuration moves from $M_0$ toward the lower-workload
configurations, the average inference throughput increases and the tail
latency decreases. The measured average FPS increases from approximately
$30.3$~FPS for $M_0$ to $50.2$~FPS for $M_1$ and $51.5$~FPS for $M_2$.
Similarly, the P95 latency decreases from approximately $367$~ms for $M_0$
to $202$~ms for $M_1$ and $147$~ms for $M_2$.

However, the improvement in computational performance is accompanied by a
reduction in the number of detections generated during the observation
period. The average number of detections decreases from approximately
$8.6$ for $M_0$ to $3.7$ for $M_1$ and $2.5$ for $M_2$. The average detection
confidence also decreases moderately, from approximately $0.575$ for $M_0$
to $0.508$ for $M_1$ and $0.485$ for $M_2$. Thus, more aggressive workload
reduction improves throughput and latency, but also reduces the temporal
frequency of detection and can affect detection confidence.

These observations show that the three operating configurations represent
different workload--performance trade-offs rather than uniformly improving
operating points. In particular, $M_1$ provides a substantial throughput
improvement over $M_0$ while maintaining a moderate reduction in detection
confidence, whereas the transition from $M_1$ to $M_2$ provides a further
reduction in P95 latency but only a small increase in average FPS and a
further reduction in detection frequency. This demonstrates that selecting
an operating point involves balancing computational efficiency against
detection availability and performance.

The results therefore motivate adaptive workload control rather than the use
of a single fixed configuration. Under changing device conditions, a
controller can select an appropriate operating point according to the
current computational and inference state. The Phase~1 measurements
therefore provide the empirical motivation for the adaptive-control stages
that follow.

\begin{figure}[t]
    \centering
    \includegraphics[width=\linewidth]{logs_rpi5_mode_summary.pdf}
    \caption{Phase~1 mode characterization on the Raspberry Pi~5. The
    lower-workload configurations increase average FPS and reduce P95
    latency, but also reduce the average number of detections and moderately
    decrease average detection confidence.}
    \label{fig:phase1_mode_summary}
\end{figure}

The discrete configurations are therefore used only for empirical
characterization and as an interpretable heuristic baseline. The final
learned controller instead operates over the continuous action parameters
$(imgsz_t,p_{skip,t})$, allowing it to select intermediate operating points
rather than being restricted to the three discrete configurations evaluated
in Phase~1. The observed workload--performance trade-off provides the
motivation for learning such a continuous adaptation policy.


\subsection{Phase 2 \& 3: Deterministic Adaptive Controller}

Phase~2 introduces closed-loop telemetry by continuously monitoring
$U_{cpu}$, $U_{ram}$, and $T_{core}$ during inference. Phase~3 connects
this telemetry to a deterministic rule-based controller that selects an
inference mode at each decision step.

The controller maps the current resource utilization directly to one of the
three modes defined in Phase~1. The decision is based on CPU and RAM
utilization as follows:

\begin{algorithm}[H]
\caption{Deterministic Adaptive Controller}
\label{alg:heuristic}
\begin{algorithmic}[1]
\Require $U_{cpu}$, $U_{ram}$
\Ensure $m_{next} \in \{M_0,M_1,M_2\}$

\If{$U_{cpu}>80.0$ \textbf{or} $U_{ram}>85.0$}
    \State $m_{next}\gets M_2$ \Comment{Severe workload reduction}
\ElsIf{$U_{cpu}>50.0$}
    \State $m_{next}\gets M_1$ \Comment{Moderate workload reduction}
\Else
    \State $m_{next}\gets M_0$ \Comment{Full-resolution inference}
\EndIf
\end{algorithmic}
\end{algorithm}

The thresholds were selected from the operating ranges observed during
Phase~0 and Phase~1. Since sustained inference produced CPU utilization
within approximately $45$--$90\%$, $50\%$ was used to trigger moderate
workload reduction, while $80\%$ was selected for severe load conditions.
The $85\%$ RAM threshold provides an additional safeguard against high
memory utilization. The controller is stateless, with each decision made
independently of the previous mode; consequently, no hysteresis or recovery
timer is applied.

This deterministic controller demonstrates that adaptive workload control
can be achieved using simple and auditable rules. However, its fixed
thresholds restrict the controller to three predefined operating points.
The subsequent learned controller therefore replaces these discrete rules
with a continuous policy over $(imgsz_t,p_{skip,t})$, enabling adaptation
across a broader range of workload conditions.

\subsection{Phase 4: Formulation as a Markov Decision Process}
To develop a predictive and optimized controller, the dynamic inference problem is formulated as a Markov Decision Process (MDP) defined by the tuple $\langle \mathcal{S}, \mathcal{A}, \mathcal{P}, \mathcal{R}, \gamma \rangle$.

\textbf{State Space ($\mathcal{S}$):} 
A continuous vector representing current physical telemetry and performance:
\begin{equation}
s_t = [U_{gpu}, U_{vram}, U_{cpu}, T_{gpu}, FPS_{avg}, M_{prev}]
\end{equation}

\textbf{Action Space ($\mathcal{A}$):}
The discrete set of modes $\mathcal{A} = \{M_0, M_1, M_2\}$.

\textbf{Reward Function ($\mathcal{R}$):}
The reward function encodes the engineering objectives: maximize throughput while maintaining thermal safety.
\begin{equation}
R_t = \alpha \cdot FPS_t - \beta \cdot P_{temp}(T_{gpu}) - \delta \cdot P_{switch}(M_t, M_{prev})
\end{equation}
Where:
- $FPS_t$ rewards the maintenance of high throughput.
- $P_{temp}$ represents a non-linear penalty function for temperatures approaching the throttling limit.
- $P_{switch}$ applies a minor penalty for mode switching to discourage jitter.
- $\alpha, \beta, \delta$ are weighting hyperparameters.

The temperature penalty is modeled exponentially:
\begin{equation}
P_{temp}(T) = 
\begin{cases} 
0 & \text{if } T \leq 65^\circ\text{C} \\
e^{\lambda (T - 65)} - 1 & \text{if } 65^\circ\text{C} < T \leq 85^\circ\text{C} \\
1000 & \text{if } T > 85^\circ\text{C} \text{ (Failure)}
\end{cases}
\end{equation}

In Phase 4, a massive offline dataset of state-action-reward tuples is generated by running the system using a combination of random exploration and the heuristic controller acting as a behavioral policy.

\subsection{Phase 5: Learning-Based Controller (DQN-lite)}
With the MDP formalized and data collected, Phase 5 trains the Learning-Based Controller. Because edge devices possess limited computational resources, running a massive deep reinforcement learning agent concurrently with YOLOv8 is impractical. Therefore, a lightweight variant of the Contextual Bandit and Deep Q-Network, referred to as DQN-lite, is utilized.

The agent's neural network comprises three fully connected layers (sizes 32, 32, 16) with ReLU activations. It accepts the 6-dimensional state vector $s_t$ as input and outputs a 3-dimensional vector of Q-values $Q(s_t, a; \theta)$ for each potential action.

The network is trained to minimize the Temporal Difference (TD) error:
\begin{equation}
L(\theta) = \mathbb{E} \left[ \left( R_{t+1} + \gamma \max_{a'} Q(s_{t+1}, a'; \theta^-) - Q(s_t, a_t; \theta) \right)^2 \right]
\end{equation}

During real-time deployment, the Controller selects the action maximizing the predicted Q-value: $a_t = \arg\max_{a} Q(s_t, a; \theta)$. The minimal size of the neural network allows a forward pass to execute in under $1$ ms on the CPU, introducing negligible latency.



\subsection{Phase 6 \& 7: Optimization and Online Adaptation}
Phase 6 introduces software optimizations to further enhance system efficiency. The YOLOv8 model is configured to operate with FP16 precision using NVIDIA TensorRT. TensorRT performs layer fusion and automatic kernel tuning optimized for the target graphics card. This significantly reduces the thermal output of the model, enabling Mode $M_0$ to generate less heat per processed frame.

Phase 7 implements online adaptation. The RL agent periodically updates its policy using a small replay buffer of recent experiences. This continuous, lightweight fine-tuning allows the controller to dynamically recalibrate its behavior in response to changes in ambient temperature or hardware degradation over time.

\section{Experimental Results}
\label{sec:experiments}
Comprehensive evaluations were conducted to validate the efficacy of the Learned Adaptive Inference System against static and heuristic baselines.

\subsection{Experimental Setup}
\textbf{Hardware and Software:} All experiments reported in the baseline and
controlled-adaptation stages were conducted on a Raspberry Pi~5, the target
resource-constrained edge platform of this study. The inference pipeline uses
Python, OpenCV for video acquisition, Ultralytics YOLOv8n for object detection,
and \texttt{psutil} for hardware telemetry. The baseline uses a fixed
$640 \times 640$ input and processes every decoded frame. The
controlled-adaptation experiment evaluates the three configurations defined in
Phase~1: $M_0=(640,0)$, $M_1=(480,1)$, and $M_2=(320,2)$. GPU telemetry is not
used as a control variable, and no desktop system, NVIDIA Jetson emulator, or
TensorRT optimization is used in these experiments.

\textbf{Dataset and Workload:} The workload is a locally stored MP4 road-scene
video containing pothole imagery, processed sequentially through OpenCV. The
baseline and controlled-adaptation scripts each use a recorded MP4 input
configured in their respective run scripts. The controlled-adaptation pipeline
uses the YOLOv8 weights stored in
\texttt{BaseModelTraining/test/best.pt}; the baseline pipeline uses the
baseline weights configured in its run script. The experiments therefore
measure sustained inference on a recorded road-scene workload rather than
online camera-network transmission or a fixed 15,000-frame synthetic test.

\textbf{Evaluation Metrics:} The runtime scripts record the following
measurements for each processed frame:
\begin{itemize}
    \item \textbf{Inference FPS and latency:} The smoothed inference frame rate
    and elapsed inference latency in milliseconds.
    \item \textbf{Resource telemetry:} CPU utilization, RAM utilization, and
    core temperature sampled asynchronously by the system monitor.
    \item \textbf{Detection output:} The number of detections, mean detection
    confidence, and the maximum-confidence proxy recorded by the runtime logger.
    \item \textbf{Frame processing:} For the controlled experiment, source
    frames seen, processed frames, intentional skips, intentional drop rate,
    and total frame-drop rate.
\end{itemize}

The runtime logs do not contain ground-truth annotations and therefore do not
directly measure mAP. Accuracy claims based on mAP or class-level validation
are reported separately only when obtained from the dedicated evaluation
pipeline; they are not treated as measurements of the runtime setup described
here.

\subsection{Baseline vs. Heuristic Performance}
The static baseline was evaluated on the Raspberry Pi~5 with a fixed
$640 \times 640$ input and no intentional frame skipping. Across the recorded
run, the baseline log contains 692 inference rows, with a mean inference rate
of $3.31$~FPS and mean latency of $301.80$~ms. The logged core temperature
ranged from $49.6^\circ$C to $60.6^\circ$C, while the final cumulative frame
drop rate was $87.04\%$. These measurements show that the main limitation of
the tested static configuration was insufficient processing throughput rather
than an observed thermal-throttle event.

The deterministic controller was evaluated in a separate Raspberry Pi~5 run
using the CPU- and RAM-based rules described in Phase~2 and Phase~3. Across
402 logged inference rows, it achieved a mean inference rate of $5.35$~FPS
and mean latency of $194.60$~ms. The final frame-drop rate decreased to
$41.91\%$, and the measured core temperature remained between $51.2^\circ$C
and $60.0^\circ$C. The controller switched between $M_0$ and $M_1$ 86 times
in this run; $M_2$ was not selected in the recorded Raspberry Pi log.

The comparison indicates that workload adaptation improved effective
processing throughput and reduced frame loss under the tested conditions.
However, the deterministic controller remains reactive and stateless: each
decision is made from the current CPU and RAM measurements, without a model of
future thermal or workload trends. This makes it interpretable and easy to
audit, but can also cause repeated mode changes and motivates the learned
controller developed in the subsequent phases. Because the baseline and
controller measurements come from separate recorded runs, the values are
interpreted as runtime characterization rather than a paired statistical
comparison.

\subsection{Learning-Based Controller Performance}
Phase~4 collected the offline transition data required to train the learned
controller. The Raspberry Pi~5 collection run produced 3,033 state--action--
reward transitions. The logged rewards ranged from $-12.72$ to $907.97$, with
an average of $683.69$. The exploration controller sampled continuous raw
actions in $[-1,1]^2$, which were mapped to image sizes from 256 to 640 pixels
and skip probabilities from 0 to 1.

Phase~5 trains a $5 \rightarrow 128 \rightarrow 128 \rightarrow 2$
Continuous Actor using PyTorch. The training script retains the top 30\% of
transitions by reward, standardizes the five state variables using saved mean
and standard-deviation arrays, and minimizes mean-squared error between the
actor output and the selected raw actions for 150 epochs. This is an offline
reward-filtered behavioral-cloning procedure; it is not a deployed DQN and it
does not perform online exploration or policy updates during inference.

The available Phase~5 artifacts include the state-normalization arrays and a
verification plot generated by gradually varying a simulated temperature
state. An exported policy file and a hardware-in-the-loop evaluation run are
not present in the current repository. Consequently, the present results
establish that the transition dataset and offline training workflow are
available, but do not support claims that the learned policy has achieved a
specific FPS, temperature, frame-drop rate, or accuracy improvement over the
deterministic controller.

\subsection{Impact of TensorRT Optimization}
TensorRT FP16 optimization is described as a possible deployment optimization,
but it is not part of the Raspberry Pi~5 baseline, controlled-adaptation, or
offline-controller results reported in this manuscript. Accordingly, no
TensorRT FPS, thermal, or accuracy improvement is claimed here. A controlled
comparison of the learned policy with and without an appropriate inference
backend remains future work.

\subsection{System Performance Evaluation Across Platforms}
\begin{figure*}[t!]
    \centering
    \includegraphics[width=0.85\textwidth]{runs_comparison_summary.png}
    \caption{Summary of runtime comparisons highlighting performance stability across different evaluation runs.}
    \label{fig:runs_comparison}
\end{figure*}

To ensure the robustness of the adaptive framework, the system's performance was evaluated across varying hardware profiles. A comprehensive comparison of key system parameters—including GPU utilization, CPU load, and thermal stability—was conducted.

\begin{figure}[H]
    \centering
    \includegraphics[width=\linewidth]{sysparams_bar_comparison.png}
    \caption{Bar chart comparison of critical system parameters across static, heuristic, and learned configurations.}
    \label{fig:sysparams_bar}
\end{figure}

As shown in Fig. \ref{fig:runs_comparison} and Fig. \ref{fig:sysparams_bar}, the learned controller demonstrated a superior ability to balance multidimensional constraints. While static pipelines maximized FPS at the expense of catastrophic thermal failure, and heuristic controllers prioritized thermal safety at the cost of erratic FPS drops, the DQN-lite agent found an optimal operating frontier. 

\begin{figure}[H]
    \centering
    \includegraphics[width=0.9\linewidth]{sysparams_radar.png}
    \caption{Radar chart providing a multidimensional view of system performance metrics, including thermal safety, FPS, and resource utilization.}
    \label{fig:sysparams_radar}
\end{figure}

The radar chart (Fig. \ref{fig:sysparams_radar}) further visualizes this balance, illustrating that the learned approach yields the most symmetrical and stable resource utilization footprint among all evaluated methods.

\subsection{Model Accuracy and Fidelity Trade-offs}
\begin{figure*}[t!]
    \centering
    \includegraphics[width=0.85\textwidth]{accuracy_comparison_bar.png}
    \caption{Bar chart comparing detection accuracy across different inference modes (M0, M1, M2).}
    \label{fig:accuracy_comparison}
\end{figure*}

A critical concern in dynamic inference is the potential loss of detection fidelity when operating in degraded modes. The impact of resolution scaling on the YOLOv8n model's accuracy was thoroughly quantified.

\begin{figure}[H]
    \centering
    \includegraphics[width=\linewidth]{per_image_f1_scatter.png}
    \caption{Scatter plot illustrating per-image F1 scores, revealing the variance in detection reliability under different visual conditions.}
    \label{fig:f1_scatter}
\end{figure}

The experimental data indicates that transitioning from $640 \times 640$ ($M_0$) to $480 \times 480$ ($M_1$) incurs a negligible drop in Mean Average Precision (mAP), typically around 2-3\%. However, dropping to $320 \times 320$ ($M_2$) results in a more pronounced degradation, as visualized in the per-image F1 scores (Fig. \ref{fig:f1_scatter}). 

\begin{figure}[H]
    \centering
    \includegraphics[width=\linewidth]{aggregate_metrics.png}
    \caption{Aggregate accuracy metrics over time, demonstrating the overall robustness of the adaptive inference approach.}
    \label{fig:aggregate_metrics}
\end{figure}

The learning-based controller's ability to maximize time spent in $M_0$ and $M_1$, while actively avoiding $M_2$, is fundamental to its success in preserving overall system accuracy over sustained periods (Fig. \ref{fig:aggregate_metrics}).

\subsection{Controller Adaptation Dynamics and Policy Analysis}
\begin{figure*}[t!]
    \centering
    \includegraphics[width=0.85\textwidth]{sysparams_timeseries.png}
    \caption{Time-series analysis of system parameters, showcasing the precise moments the learned controller triggers preemptive mode shifts to avert thermal crises.}
    \label{fig:sysparams_timeseries}
\end{figure*}

The internal behavior of the learned policy provides deeper insights into its efficacy. An analysis of the agent's action distribution reveals its proactive nature, as seen in the system parameter time-series (Fig. \ref{fig:sysparams_timeseries}). 

\begin{figure}[H]
    \centering
    \includegraphics[width=\linewidth]{rl_imgsz_distribution.png}
    \caption{Distribution of inference modes selected by the RL agent during sustained operation.}
    \label{fig:imgsz_dist}
\end{figure}

The mode distribution (Fig. \ref{fig:imgsz_dist}) highlights a strong preference for $M_1$ as a sustainable operating point, confirming the hypothesis that preemptive moderate throttling is vastly superior to reactive aggressive throttling.

\begin{figure}[H]
    \centering
    \includegraphics[width=\linewidth]{adaptation_curve.png}
    \caption{Adaptation curve illustrating the RL agent's learning progression and the stabilization of the reward signal.}
    \label{fig:adaptation_curve}
\end{figure}

Finally, the adaptation curve (Fig. \ref{fig:adaptation_curve}) demonstrates that the agent quickly learns to avoid the severe penalties associated with thermal throttling, stabilizing its reward signal efficiently during the training phase.

\section{Discussion}
\label{sec:discussion}
The experimental results validate the core hypothesis: edge inference must be treated as a dynamic, hardware-aware process. 

\subsection{The Value of Anticipation}
The fundamental advantage of the learning-based approach over rule-based heuristics is its capacity for anticipation. Thermal systems possess significant inertia; when a temperature sensor registers $80^\circ$C, the heatsink is already saturated, and temperatures will continue to rise even if the workload is immediately reduced. A rule-based system reacts only after this momentum is established, necessitating aggressive fallback measures.

Conversely, the Reinforcement Learning agent internalizes the consequences of thermal momentum through its updated Q-values. It recognizes that preemptively shifting to $M_1$ yields a higher long-term reward than sustaining $M_0$ until an emergency $M_2$ shutdown is unavoidable. By anticipating delayed consequences, the learned controller avoids severe performance penalties and maintains stable operation.

\subsection{Interpretation of Baseline and Deterministic Control Results}

The baseline and deterministic-controller measurements justify the use of
adaptive workload control at the application layer. The fixed $640 \times 640$
pipeline processed frames slowly and discarded most source frames because each
frame required a substantial inference time. The deterministic controller
reduced this mismatch by selecting lower-workload configurations when resource
utilization increased, improving the measured throughput and reducing the
frame-drop rate while keeping the observed core temperature below $60^\circ$C.

These results should not be interpreted as evidence that the controller
eliminates all frame loss or guarantees hard real-time behavior. The remaining
frame drops and repeated mode changes demonstrate the limitation of a
stateless, threshold-based policy. This limitation provides the engineering
justification for the later learned controller, whose objective is to select
workload parameters using the longer-term consequences represented in the
state--action--reward data.

\subsection{Interpretation of Offline Controller Development}

The Phase~4 and Phase~5 results establish the data and training stages of the
learned-control pipeline, but they do not constitute an end-to-end deployment
evaluation. The transition data provides coverage of continuous workload
actions and their observed rewards, while reward filtering focuses training on
actions associated with better measured operating outcomes. Because the
available artifacts do not include a completed hardware-in-the-loop policy
run, the learned controller should not yet be credited with the performance
of the deterministic controller or with a measured accuracy improvement.

This distinction is important for interpreting the contribution: offline
behavioral cloning removes exploratory policy updates from deployment, but it
does not by itself prove that the learned policy generalizes to unseen thermal
conditions. A subsequent evaluation should load the exported policy on the
Raspberry Pi~5, replay the same workload, and report paired FPS, latency,
temperature, frame-drop, and detection metrics against the static and
deterministic baselines.

\subsection{Local Adaptive Inference Versus Edge-Cloud Offloading}

Edge-cloud task offloading provides an alternative mechanism for addressing
computational limitations at resource-constrained devices. By distributing
tasks between the local device, edge servers, and cloud infrastructure,
offloading can provide access to additional computational resources while
jointly considering communication and processing constraints
\cite{b25}. This approach is particularly relevant for applications in which
network connectivity and bounded communication latency can be reliably
provided.

The present work instead adopts a connectivity-independent architecture in
which both inference and control remain on the edge device. This design choice
is motivated by the deployment objective of maintaining local road-surface
monitoring without requiring continuous communication with an external
computing infrastructure. Offloading would introduce an additional
communication path whose latency and availability depend on network
conditions, in addition to the computational cost of transmitting video data
or inference inputs. These factors are outside the scope of the current
experimental evaluation.

This distinction does not imply that edge-cloud offloading is unsuitable for
adaptive road-monitoring systems. Rather, the two approaches address
different deployment assumptions. The proposed method reduces computational
demand locally through adaptive resolution and frame-processing frequency,
whereas an IoT-edge-cloud architecture can redistribute computation when
suitable communication and remote computational resources are available.
Integrating the proposed local controller with deterministic task offloading
could therefore form a complementary hybrid architecture in future work,
where local adaptation provides immediate workload control and offloading is
invoked when network and remote-computing conditions satisfy the required
service constraints.

\subsection{Generalization to Other Models}
While this study utilized YOLOv8n, the Learned Adaptive Inference System is inherently model-agnostic. The defined Action Space modulates universal pipeline parameters—input resolution and frame-skip intervals—rather than internal model weights. Consequently, the controller can be seamlessly integrated with ResNet classification pipelines, Mask R-CNN segmentation models, or transformer-based architectures, provided the Action Space is reprofiled for the specific model's computational cost.

\subsection{Limitations}
A primary limitation of the current implementation is that the Action Space does not account for the visual complexity or context of the video stream. For instance, if a vehicle is stationary at a traffic light, dropping 90\% of frames would have negligible impact on hazard tracking while saving substantial power. 

Future iterations could incorporate content-aware metrics, such as "scene entropy" or "optical flow magnitude," into the state vector. This would enable the controller to modulate workloads based on visual scene dynamics in addition to hardware constraints, further optimizing power efficiency during idle periods.

\section{Conclusion}
\label{sec:conclusion}
Deploying sophisticated computer vision models on resource-constrained edge devices requires a paradigm shift away from static inference pipelines. This paper detailed the development of a Learned Adaptive Inference System that dynamically adjusts the operational parameters of a YOLOv8 pothole detection model based on real-time hardware telemetry.

The results demonstrated that while rule-based heuristic controllers can prevent thermal throttling, their reactive nature leads to oscillatory performance and unnecessary frame drops. By formulating the problem as a Markov Decision Process and employing a DQN-lite agent, the system achieved anticipatory thermal management. The learned controller maintained stable temperatures while maximizing throughput, ensuring continuous and reliable hazard detection for autonomous navigation. 

Furthermore, combining the adaptive pipeline with TensorRT optimization yielded compounding benefits, proving that software-level adaptivity and hardware-level optimization are highly complementary. Future research will focus on integrating content-aware features and exploring federated learning techniques to enable cross-device policy transfer across heterogeneous edge platforms.

\begin{thebibliography}{99}

\bibitem{b1}
R. Kothai and N. Prabakaran, ``Self Attention GAN and SWIN Transformer-Based Pothole Detection With Trust Region-Based LSM and Hough Line Transform for 2D to 3D Conversion,'' \textit{IEEE Access}, vol. 13, pp. 137695--137708, Aug. 2025, doi: 10.1109/ACCESS.2025.3595190.

\bibitem{b2}
S. Kong, Q. Meng, X. Li, Z. Wang, X. Liu, and B. Li, ``Road Perception for Autonomous Driving: Pothole Detection in Complex Environments Based on Improved YOLOv8,'' \textit{IEEE Access}, vol. 13, pp. 80914--80929, May 2025, doi: 10.1109/ACCESS.2025.3566944.

\bibitem{b3}
G. Shelke et al., ``Vision-Based Pothole Detection and Reporting System for Smart Cities (Leveraging YOLOv8-Seg and Predictive Maintenance Models),'' in \textit{Proc. IEEE Pune Section Int. Conf. (PuneCon)}, 2025, doi: 10.1109/PuneCon67554.2025.11379854.

\bibitem{b4}
K. Agarwal, S. Bhambri, V. K. Sridharan, N. Mohammed, H. Kohli, and J. A. Kapoor, ``Performance Evaluation of different Machine Learning Techniques for Pothole Detection,'' in \textit{Proc. IEEE Int. Conf. Contemporary Computing and Communications (InC4)}, 2025, doi: 10.1109/InC465408.2025.11256464.

\bibitem{b5}
M. Yurdakul and S. Tasdemir, ``An Enhanced YOLOv8 Model for Real-Time and Accurate Pothole Detection and Measurement,'' arXiv preprint arXiv:2505.04207, May 2025.

\bibitem{b6}
Y. Hu et al., ``Lightweight Deep Learning for Real-Time Road Distress Detection on Mobile Devices,'' \textit{Nature Communications}, vol. 16, Art. no. 4212, May 2025.

\bibitem{b7}
N. Yang, S. Chen, H. Zhang, and R. A. Berry, ``Beyond the Edge: An Advanced Exploration of Reinforcement Learning for Mobile Edge Computing, Its Applications, and Future Research Trajectories,'' \textit{IEEE Communications Surveys \& Tutorials}, vol. 27, no. 1, pp. 546--594, 2025, doi: 10.1109/COMST.2024.3405075.

\bibitem{b8}
K. Ye et al., ``Energy-Efficient Inference With Software--Hardware Co-Design for Sustainable Artificial Intelligence of Things,'' \textit{IEEE Internet of Things Journal}, vol. 11, no. 24, pp. 39170--39182, Dec. 2024, doi: 10.1109/JIOT.2024.3482288.

\bibitem{b9}
X. Chen, J. Cao, Y. Sahni, S. Jiang, and Z. Liang, ``Dynamic Task Offloading in Edge Computing Based on Dependency-Aware Reinforcement Learning,'' \textit{IEEE Transactions on Cloud Computing}, vol. 12, no. 2, pp. 594--608, Apr.--Jun. 2024, doi: 10.1109/TCC.2024.3381646.

\bibitem{b10}
G. P. B. et al., ``Accelerating Native Inference Model Performance in Edge Devices using TensorRT,'' in \textit{Proc. IEEE Recent Advances in Intelligent Computational Systems (RAICS)}, 2024, doi: 10.1109/RAICS61201.2024.10690032.

\bibitem{b11}
``Enhancing Infrastructure Monitoring: Pothole Detection in Road Images Using YOLOv8 and Open Datasets,'' in \textit{Proc. IEEE 6th Int. Conf. on Cybernetics and Intelligent System (ICORIS)}, 2024, doi: 10.1109/ICORIS63540.2024.10903697.

\bibitem{b12}
``Pothole Detection and Filling System Using Image Processing and Machine Learning,'' in \textit{Proc. Int. Conf. Smart Systems for Applications in Electrical Sciences (ICSSES)}, 2024, doi: 10.1109/ICSSES62373.2024.10561417.

\bibitem{b13}
S. D. RS et al., ``Pothole Detection and Instance Segmentation Using YOLO V8,'' in \textit{Proc. IEEE Int. Conf. IoT Based Control Networks and Intelligent Systems (ICICNIS)}, 2024, doi: 10.1109/ICICNIS64247.2024.10823139.

\bibitem{b14}
R. C. C. M. Santos, M. Coelho, and R. Oliveira, ``Real-Time Object Detection Performance Analysis Using YOLOv7 on Edge Devices,'' \textit{IEEE Latin America Transactions}, vol. 22, no. 10, pp. 799--805, Oct. 2024, doi: 10.1109/TLA.2024.10705971.

\bibitem{b15}
R. M. et al., ``An Annotated Street View Image Dataset for Automated Road Damage Detection (SVRDD),'' \textit{Scientific Data}, vol. 11, Art. no. 407, Apr. 2024.

\bibitem{b16}
A. R. Khouas et al., ``Training Machine Learning Models at the Edge: A Survey,'' arXiv preprint, Jul. 2024.

\bibitem{b17}
``AdaDet: An Adaptive Object Detection System Based on Early-Exit Neural Networks,'' \textit{IEEE Transactions on Cognitive and Developmental Systems}, vol. 16, no. 1, pp. 332--345, Feb. 2024, doi: 10.1109/TCDS.2023.3274214.

\bibitem{b18}
X. Zhang, Y. Teng, N. Wang, B. Sun, and G. Hu, ``Accelerating Deep Neural Network Tasks Through Edge-Device Adaptive Inference,'' in \textit{Proc. IEEE 34th Annual Int. Symp. Personal, Indoor and Mobile Radio Communications (PIMRC)}, 2023, pp. 1--6, doi: 10.1109/PIMRC56721.2023.10293996.

\bibitem{b19}
S. Maity, R. Roy, A. Majumder, S. Dey, and A. R. Hota, ``Future-Aware Dynamic Thermal Management in CPU-GPU Embedded Platforms,'' in \textit{Proc. IEEE Real-Time Systems Symposium (RTSS)}, 2022, doi: 10.1109/RTSS55097.2022.00041.

\bibitem{b20}
G. Verma, Y. Gupta, A. M. Malik, and B. Chapman, ``Performance Evaluation of Deep Learning Compilers for Edge Inference,'' in \textit{Proc. IEEE Int. Parallel and Distributed Processing Symposium Workshops (IPDPSW)}, 2021, pp. 858--865, doi: 10.1109/IPDPSW52791.2021.00128.

\bibitem{b21}
X. Gao, Y. Zhao, L. Dudziak, R. Mullins, and C. Xu, ``Dynamic Channel Pruning: Feature Boosting and Suppression,'' in \textit{Proc. Int. Conf. Learning Representations (ICLR)}, 2019.

\bibitem{b22}
B. Taylor, V. S. Marco, W. Wolff, Y. Elkhatib, and Z. Wang, ``Adaptive Deep Learning Model Selection on Embedded Systems,'' in \textit{Proc. 19th ACM SIGPLAN/SIGOPS Int. Conf. Virtual Execution Environments}, 2018, pp. 31--43.

\bibitem{b23}
S. Teerapittayanon, B. McDanel, and H. T. Kung, ``BranchyNet: Fast Inference via Early Exiting from Deep Neural Networks,'' in \textit{Proc. 23rd Int. Conf. Pattern Recognition (ICPR)}, 2016, pp. 2464--2469.

\bibitem{b24}
M. Mutapcic and S. Boyd, ``Cutting-Set Methods for Robust Convex Optimization with Pessimizing Oracles,'' \textit{Optimization Methods and Software}, vol. 24, no. 3, pp. 381--406, 2009.

\bibitem{b25}
K. Aghababaiyan, B. Coll-Perales, and J. Gozalvez,
``Deterministic Task Offloading and Resource Allocation in the
IoT-Edge-Cloud Continuum,'' in \textit{Proc. IEEE 101st Vehicular
Technology Conference (VTC2025-Spring)}, 2025, pp. 1--6,
doi: 10.1109/VTC2025-Spring65109.2025.11174312.

\end{thebibliography}

\begin{IEEEbiography}[{\includegraphics[width=1in,height=1.25in,clip,keepaspectratio]{placeholder_a2.png}}]{Vishal K. Bhat}
is pursuing his Bachelor's degree in Computer Science at RV College of Engineering, Bengaluru, India. His research interests encompass computer vision, edge computing, and IoT systems. In this project, he led the development of the YOLOv8 detection pipeline and implemented the baseline profiling and TensorRT optimizations.
\end{IEEEbiography}

\begin{IEEEbiography}[{\includegraphics[width=1in,height=1.25in,clip,keepaspectratio]{placeholder_a3.png}}]{Vaibhav P. Rodappanavar}
is pursuing his Bachelor's degree in Computer Science at RV College of Engineering, Bengaluru, India. His research interests include machine learning, embedded systems, and adaptive inference algorithms. He contributed to the development of the reinforcement learning framework and the hardware abstraction layer for real-time system monitoring.
\end{IEEEbiography}

\begin{IEEEbiography}[{\includegraphics[width=1in,height=1.25in,clip,keepaspectratio]{placeholder_a4.png}}]{Naren Sridharan}
is pursuing his Bachelor's degree at RV College of Engineering, Bengaluru, India. His expertise includes secure edge deployments and system integration. He contributed to the experimental design, dataset management, and ensuring robust system recovery during thermal overload testing.
\end{IEEEbiography}

\begin{IEEEbiography}[{\includegraphics[width=1in,height=1.25in,clip,keepaspectratio]{placeholder_a5.png}}]{Hardik Paneru}
is pursuing his Bachelor's degree at RV College of Engineering, Bengaluru, India. His research explores deep learning optimization and embedded software engineering. He formulated the Markov Decision Process (MDP) and contributed to the architecture utilized in the learned controller.
\end{IEEEbiography}

\begin{IEEEbiography}[{\includegraphics[width=1in,height=1.25in,clip,keepaspectratio]{placeholder_a1.png}}]{Rajeswara Rao K V S}
is the Head of the Department of Industrial Engineering \& Management at RV College of Engineering, Bengaluru, India. His research interests span industrial systems, process optimization, and the integration of artificial intelligence in engineering applications. He provided critical leadership, theoretical guidance, and project oversight for the adaptive edge inference system. ORCID: https://orcid.org/0000-0003-1592-5578.
\end{IEEEbiography}

\EOD
\end{document}
